\section{More Insight into Christian Initiation}

This is a continuation of posts on Clement's Journal, found (in obscurity) in the Ante Nicene Fathers. Peter has been initiating Clement, so we see the ``inner workings'' of this process, which really amounts to a sort of personal spiritual apprenticeship, which is governed by affectionate command and proud obedience, more than a set rule. They have been debating the doctrine of ``Genesis'' (Fate) with an old man who is an honest hermit (and will turn out to be Clement's long lost father.

The following may sound like ``the Fall'', but what is being said (if one can hear) is that the very elements were tainted\footnote{\url{https://en.wikipedia.org/wiki/One\_Power}} by the spiritual worship offered to demons:

\begin{quotex}
But by the freedom of the will, every man, while he is unbelieving in regard to things to come, runs into evils by doing evil deeds. And these are the things in the world which seem to be done contrary to order, which owe their existence to mans rebellion against Torah. Therefore the purpose of YHWH is for man to worship Him, which He granted to men in the beginning. They were walking in the good way of life in order to enjoy incorruptible good things. But when they sinned, they gave birth to evil by sin. And to every good thing evil is joined as by a certain covenant of alliance on the part of sin, since indeed the earth has been polluted with human blood, and altars have been lighted to demons, and they have polluted the very air by the filthy smoke of zebahim; and so at length the elements, being first corrupted, have handed over to men the fault of their corruption, as roots nourish the branches and the fruit... 
\end{quotex}

Thus, we see, by a divine circle, being brought back to the opening argument of Simon Magus that Christ's words about ``divison'' and ``houses falling'' were nonsense, that G-D sent Christ into the world to cleave good and evil asunder -- Christ is God's ``sword''.

Peter offers another good argument as to why unbelief ultimately needs a ``First Principle'' from Revelation in order to truly begin to deal with Reality:

\begin{quotex}
Hence I think it ridiculous when men judge of the power of Elohim in natural ways, and think that this is possible and that impossible to Him, or this greater and that less,while they are ignorant of everything; who, being unrighteous, judge the righteous YHWH; unskilled, judge the contriver; corrupt, judge the incorruptible; created, judge the Creator... 
\end{quotex}

This makes a deal of sense --- we ordinarily let men ``speak'' for themselves, why not God? Is He worthy of less privilege than a mere man? Reading on, we are told that Jesus' words are ``simple, plain, and brief'', hence, Platonism for the masses, plus something. Clement tries to summarize:

\begin{quotex}
YHWH by His Son created the world as a double house,separated by the interposition of this firmament, which is called heaven; and appointed malachim to dwell in the higher, and a multitude of men to be born in this visible world, from amongst whom He might choose friends for His Son, with whom He might rejoice, and who might be prepared for Him as a beloved bride for a bride groom. But even till the time of the marriage, which is the revelation of the world to come, He has appointed a certain power,to choose out and watch over the good ones of those who are born in this world, and to preserve them for His Son, who is without sin, set apart in a certain place of the world, in which there are already some who are there being prepared, as I said, as a bride adorned for the coming of the bridegroom. For the prince of this world and of the present age is like an adulterer who corrupts and violates the minds of men, and, seducing them from the love of the true bridegroom,allures them to strange lovers... 
\end{quotex}


There is a ``certain power'' which preserves those in this world who inherit Life --- the ``leaven of God'' we may call it. There is also what might be called ``General Law'':

\begin{quotex}
And because He foresaw that the present world could not consist except by variety and inequality, He gave to each mind freedom of motions,according to the diversities of present things, and appointed this prince, through that princes own suggestion of those things which run contrary, that the choice of better things might depend upon the exercise of virtue... 
\end{quotex}


Baptism is explained, again:

\begin{quotex}
For He has given a Torah, thereby aiding the minds of men, that they may the more easily perceive how they ought to act with respect to everything, in what way they may escape evil, and in what way tend to future blessings; and how, being regenerate in water, they may by good works extinguish the fire of their old birth. For our first birth descends through the fire of lust,and therefore, by the divine appointment, this second birth is introduced by water, which may extinguish the nature of fire; and that the spirit enlightened by YHWH is set apart. Spirit may cast away the fear of the first birth, provided, however, it so live for the time to come... 
\end{quotex}


Thus, baptism regenerates, but not infallibly -- it is a ``sealing''\footnote{\url{http://www.latrobe.edu.au/eyeoftheheart/assets/issue4/baptism.pdf}}. The world is a gymnasium, where we learn to exercise the virtues in resisting evil.

\begin{quotex}
For Elohim has given a nature to men by which they may be taught concerning what is good and to resist evil; that is, they may learn arts and to resist vain pleasures,and to set the Torah of Elohim before them in all things. And for this end He has permitted certain contrary powers to wander up and down in the world and to strive against us, for the reasons which have been stated before, that by striving with them the palm of victory and the merit of rewards may accrue to the righteous
\end{quotex}


Here we see a justification for the liberal arts (profane and sacred).

We also see some interesting doctrines on conception and ``love'':

\begin{quotex}
From this, therefore, sometimes the result of any persons acting incontinently and being willing not so much to resist as to yield to give harbor to these impulses in themselves by their noxious breath an intemperate, ill-conditioned, and diseased progeny is begotten. For while lust is supposedly gratified but no care is taken in the copulation, undoubtedly a weak generation is affected with the defects and frailties of those demons who instigate these evil deeds. And therefore parents are responsible for their children's defects of this sort, because they have not observed the Torah of association. Though there are also more secret causes, by which spirits are made subject to these evils, which it is not to our present purpose to state, yet it behooves every one to acknowledge the Way of YHWH, that he may learn from it the observance of generation and avoid causes of impurity, that that which is begotten may be pure.
\end{quotex}


Even fear is justified:

\begin{quotex}
But if those who believe in Elohim, and who confess the judgment to come and the penalty of ageless fire, if they do not refrain from sin, it is certain that they do not believe with full faith: for if faith is certain, fear also becomes certain. But if there be any defect in faith, fear also is weakened, and then the contrary powers find opportunity of entering.
\end{quotex}


Apparently, unrestrained ``love'' leads to the river Pyriphlegethon. Also, the origin of occultism is explained:

\begin{quotex}
Therefore the astrologers, being ignorant of such mysteries,think that these things come by the courses of the heavenly bodies. Hence also, in their answers to those who go to them to consult themas to future things, they are deceived in very many instances. Nor is it to be wondered at, for they are not naviïm; but, by long practice,the authors of errors find a sort of refuge in those things by which they were deceived, and introduce certain Climacteric Periods, that they may pretend a knowledge of uncertain things. For they represent these Climacterics as times of danger, in which one sometimes is destroyed, sometimes is not destroyed, not knowing that it is not the course of the stars but the operation of demons that regulates these things.
\end{quotex}


Revolutionaries may want to consider the following:

\begin{quotex}
This, however, is in general to be noticed, that if any are evil, not so much in their mind as in their doings, and are not borne to sin under the incitement of purpose, upon them punishment is inflicted more speedily, and more in the present life; for everywhere and always Elohim renders to every one according to his deeds as He judges to be expedient. But those who purposely practice wickedness, so that they sometimes even rage against those from whom they have received benefits, and who take no thought for teshuvah their punishment He defers to the future. For these men do not, like those of whom we spoke before, deserve to end the punishment of their crimes in the present life.
\end{quotex}


Clement argues nurture against nature:

\begin{quotex}
There are likewise amongst the Bactrians in the Indian countries immense multitudes of Brahmans, who also themselves from the tradition of their ancestors and peaceful customs and laws, neither commit murder nor adultery, nor worship idols, nor have the practice of eating animal food, are never drunk, never do anything maliciously, but always fear Elohim.
\end{quotex}


So we see a recognition (there) of natural virtue and common religion (provided it is free of demonism). However, the endless ``search for Truth'' at the expense of all sense is condemned:

\begin{quotex}
If we do not hold by the principles that we have acknowledged and confessed, but if those things which have been defined are always loosened by forgetfulness, we shall seem to be weaving Penelope's web, undoing what we have done. And therefore we ought either not to acquiesce too easily, before we have diligently examined the doctrine propounded; or if we have once acquiesced, and the proposition has been agreed to, then we ought to keep by what has been once determined, that we may go on with our inquiries respecting other matters
\end{quotex}


I would like to point out that these posts demonstrate the profound comprehensiveness, as well as ``exoteric esotericism'' of early Christian religion -- it's ability to hold its own philosophically with deep insights into the Kali Yuga (eg., the tainting of elements, demonism), the practical wisdom of Peter (a desire to instruct and lay stable foundations), the sacraments (baptism), and also the open-mindedness and true ``Catholicism'' of that Faith (eg., praise for Bactrian Brahmas). Additionally, these posts should make clear that more than a ``precise way'', we should struggle to uproot modern foundations and lay a firm one of ``sane, well-bred, masculinity'', which can then support the superstructure of a particular ``initiation''. Thus, the thread runs from the Mediterranean, through Old Europe, to us.


\flrightit{Posted on 2012-08-31 by Logres}

\begin{center}* * *\end{center}

\begin{footnotesize}
\begin{sffamily}
\texttt{Logres on 2012-09-01 at 17:51 said:}

Worth reposting (Cologero has cited): For that thing itself, which is now called the Christian religion, used to exist and was not lacking amongst the ancients from the beginning of the human race until Christ himself came in the flesh, from which time the true religion began to be called ``Christian.'' And after his resurrection and ascension into heaven, the Apostles had begun to proclaim him, and many believed; first in Antioch (as it is written) the disciples were called ``Christians.'' Therefore I have said: `This is the Christian religion in our time,' not because it did not exist previously, but because it received this name in later times.

\url{http://www.latrobe.edu.au/eyeoftheheart/assets/issue1/Golgotha.pdf}



\hfill

\end{sffamily}
\end{footnotesize}

